\subsection{Headless Firmware Generation with STM32CubeMX CLI}

The automation of the development cycle requires the ability to generate the baseline firmware project without human intervention. While STM32CubeMX is primarily known for its graphical user interface, it provides a powerful \textit{Headless Script Mode} that allows for the programmatic generation of initialization code through a Command Line Interface (CLI).

\subsubsection{Workflow Automation and CLI Scripting}

To enable deterministic and repeatable firmware generation, I developed a scripting strategy based on STM32CubeMX command files (typically \texttt{.txt} or \texttt{.script} files). This approach bypasses the GUI, loading a Micro-controller Unit (MCU) configuration or a pre-existing project file (\texttt{.ioc}), and executing the code generation process according to predefined parameters.

A typical automation script consists of several sequential commands as shown in Listing~\ref{lst:cubemx_script}.

\begin{lstlisting}[language=bash, caption={STM32CubeMX automation script example.}, label={lst:cubemx_script}]
# cubemx_script.txt
load STM32F401VCHx
config load "/path/to/project.ioc"
project name MySTM32Project
project toolchain "STM32CubeIDE"
project path /Users/michele/Desktop/output
project generate
exit
\end{lstlisting}

The execution of the generation process is then triggered from the system terminal using the CubeMX executable with the \texttt{-q} (quiet/script) flag:

\begin{lstlisting}[language=bash, caption={CLI command for headless generation.}]
/Applications/STMicroelectronics/STM32CubeMX.app/Contents/MacOS/STM32CubeMX \
  -q /path/to/cubemx_script.txt
\end{lstlisting}

\subsubsection{Technical Challenges and Solutions}

During the implementation of this headless workflow, several environment-specific challenges were encountered, particularly concerning the interaction between the Java-based executable and the host operating system.

\begin{itemize}
    \item \textbf{Residual GUI Dependencies}: On certain operating systems (such as macOS and Linux), the STM32CubeMX executable may still attempt to initiate graphical components even when running in script mode. This can lead to execution failures in headless environments (e.g., remote servers or Docker without X11). The issue was mitigated by ensuring appropriate environment variables were set or, where possible, using the \texttt{loadboard} command with specific parameters to avoid interactive peripheral configuration dialogs.
    
    \item \textbf{Peripheral Configuration Conflicts}: When loading a specific MCU model (\texttt{load} command) versus a specific development board (\texttt{loadboard}), the CLI behaves differently regarding default peripheral states. In my workflow, I prioritized the \texttt{config load} transition to ensure that user-defined pinouts and clock trees are precisely preserved during the automated generation cycle.
\end{itemize}

By abstracting the firmware generation into a CLI-driven process, the overall framework can treat the creation of the C project as an atomic step within the larger Agentic MLOps pipeline, enabling rapid transitions between model discovery, code generation, and physical deployment.
